\chapter{碳为什么能搭出如此多的结构}

\begin{kaichang}
第一册的最后一章，我们留下了一个悬案：你的皮肤、头发、肌肉，你吃的米饭和鸡蛋，这些由成百上千个原子搭成的复杂分子，凭什么能存在？当时我们预告过，答案要交给一种特殊的原子——碳。现在，第二册开场，该破案了。

先做第一册开场那样的游戏。请你在身边找四样东西：一支铅笔（用的是笔芯，不是木头）、一块方糖、一个塑料袋，还有——你自己的手背。这四样东西的化学世界里住着同一位主角：碳。铅笔芯几乎就是纯碳；糖的骨架全是碳；塑料袋是一根根长得离谱的碳链；而你体重里除去水分，剩下的“干货”将近一半是碳搭的。

再补一个更刺激的对比：铅笔芯又黑又软，一划就掉渣；钻石又透又硬，是自然界最硬的物质。可它们都是\textbf{纯碳}，原子一模一样，一个没有多，一个也没有少。

先猜一猜：同一种原子，凭什么搭出黑的和透的、软的和硬的、能吃的和烧得着的？把你的猜测写下来，这一章结束时回来对答案。
\end{kaichang}

\section{一种原子，一座动物园}

先把能看见的事实摆出来。

含碳的物质多到离谱。到今天为止，化学家登记在册的已知物质超过两亿种，其中\textbf{绝大多数}是含碳化合物——不含碳的“无机物”加起来只是零头。而且每年都有几十万种新的含碳分子被合成出来，这个名单还在疯长。因为这类分子最早都从动植物体内提取，化学家给了它们一个老名字：\textbf{有机化合物}。这个名字的历史包袱我们后面再拆，现在先记住：有机化学，基本上就是碳的化学。

多看一眼这些物质的“性格跨度”：甲烷是气体，一点就着；汽油是液体，见水不溶；白糖是固体，入水即化；酒精挥发得飞快，还辣嗓子；橡胶能拉长又弹回；尼龙能拉成比钢丝还结实的纤维；至于钻石和石墨，干脆是同一元素的两种极端。它们的核心骨架都是碳，可性格几乎没有重样的。

回忆一下贯穿全书的五个追问：①它由什么组成？②组成部分怎样连接和排列？③什么能量和概率规则决定它能否变化？④变化有多快、怎样受控制？⑤在生命中它怎样被保存、复制、调节和选择？第一册主要回答了①和③④的一部分。而这一章，以及整个第五篇，主战场是第②问：\textbf{连接和排列}。你会看到，光是“怎么连、怎么排”这一件事，就足以解释碳的百万大军。

\section{碳的四只手}

把镜头推进到原子层。碳原子核里有 6 个质子，核外 6 个电子，按第 9 章的楼层规则排布：第一层住 2 个，第二层住 4 个。关键就在这最外层的 4 个电子。

4 是个尴尬的数字。回忆第 18 章：钠最外层 1 个电子，丢掉的代价小，所以痛快地交出去当离子；氯最外层 7 个，抢 1 个就住满，所以见电子就抢。可碳呢？丢掉 4 个电子，要连续付四笔越来越贵的“电离费”；抢来 4 个电子，又要把四个负电荷硬塞进一个原子里，排斥大得离谱。两条路都太贵，于是碳走第三条路——第 19 章讲过的\textbf{共价键}：不送不抢，和邻居共享电子对。而最外层 4 个电子，恰好允许它同时参与 4 对共享，也就是说：

\textbf{一个碳原子，最多能伸出四只“手”（形成四根共价键）。}

四只手还不是碳独有的本事——硅也是四只手，氮能伸三只，氧能伸两只。碳真正的独门绝技，是它的手\textbf{拉自己的手也特别牢}。碳—碳单键的键能大约是 \ud{350}{kJ/mol} 上下（第 17 章讲过键能：拆开一根键平均要付的能量），这个数值相当微妙：够牢，链子在体温和室温下不会自己散架；又不至于牢到拆不动，需要改造时付得起“拆迁费”。对比一下邻居们：氧—氧、氮—氮单键只有碳—碳的一半左右，因为两边原子上的孤对电子互相排斥，链子容易断；硅—硅键又太长太弱，硅链在空气里很快散伙。元素周期表第 14 族一楼到五楼住了一圈，只有碳把“自己连成长链”这件事干得又好又稳。

所以碳的世界是这么开工的：碳原子和碳原子手拉手，连成链；每只碳手一次最多拉四家，多出来的手再去抓氢、氧、氮这些“挂件”。链可以伸长，可以分叉，可以首尾相咬成环——下一节就看这些花样。

\section{链、环、枝杈：结构数量的爆炸}

先从最简单的规则玩起：每个碳出四只手，没用上的手都抓氢。两个碳连成 \chem{CH_3-CH_3}（乙烷），三个连成 \chem{CH_3-CH_2-CH_3}（丙烷），都只有一种连法。可到四个碳，岔路出现了：四个碳可以排成一条直线，也可以三个成链、第四个挂在中间那个碳上，像树干分出一根枝。五碳的连法有 3 种，六碳有 5 种……

这就是碳结构的第一个秘密：\textbf{同样的原子数目，连接方式可以不同}。而连接方式的数目，随碳数增长得极其凶残。

\begin{qiaoliang}
只含碳和氢、且碳碳之间全是单键的分子（烷烃，第 39 章细讲），设碳数为 $n$，问：一共有多少种不同的连接方式？数学家把这个数数问题彻底算清了，前几项是：

\begin{center}
\begin{tabular}{cccccccc}
\toprule
碳数 $n$ & 4 & 5 & 6 & 8 & 10 & 15 & 20 \\
\midrule
连接方式数 & 2 & 3 & 5 & 18 & 75 & 4347 & 366319 \\
\bottomrule
\end{tabular}
\end{center}

看趋势：每加一个碳，种数大致翻一倍还多——这是\textbf{指数级}的增长。30 个碳的烷烃 \chem{C_{30}H_{62}}，理论上有约 $4.1\times 10^{9}$ 种连接方式。做个数量级估算：假如你一秒钟画一种结构，一年约 $3.2\times 10^{7}$ 秒，画完所有 30 碳烷烃要 $4.1\times 10^{9}\div 3.2\times 10^{7}\approx 130$ 年。而这还只是 30 个碳、还只许单键！你体内一个蛋白质分子就有几千个碳。现在明白了吗——碳不是“能搭很多结构”，而是能搭出\textbf{多到整个宇宙都造不完}的结构。
\end{qiaoliang}

除了分叉，碳链还能首尾相连咬成环。六个碳咬成的六元环特别常见：葡萄糖分子里有一个，汽油和药物分子里也到处是。环可以单挂，可以并联，可以桥接，像自行车链条一样环环相套。链、枝、环三种基本动作自由组合，就是碳骨架的全部语法。规则只有一条：每个碳，永远四只手。

\section{碳是立体的：四面体、会转的关节和锁死的关节}

到目前为止我们都把分子画在纸面上，好像碳链是平面的。错。碳是立体的。

回忆第 20 章的电子对排斥模型：一个碳伸出四根键，这四对电子互相排斥，彼此离得越远能量越低。在平面上，四根键最远只能摆成十字，夹角 $90^{\circ}$；而在三维空间里，四根键可以指向一个\textbf{正四面体}的四个顶点，两两夹角约 $109.5^{\circ}$——这才是能量最低的排布。所以，一个连了四个邻居的碳，真实形状不是十字，而是四面体：碳坐在正中央，四只手伸向四个角。

这个细节带来两个后果深远的“关节规则”：

\begin{itemize}[itemsep=2pt]
\item \textbf{单键是会转的关节}。两个碳之间只共享一对电子时，这根键像一根圆轴，两端可以绕着轴相对转动。于是一条碳链在液体里其实不停地扭动、折叠、甩动，像一条不断变换姿势的蛇。同一根链的这些不同“姿势”不算是不同的分子——转一转就回来了。
\item \textbf{双键是锁死的关节}。两个碳共享两对电子（第 19 章讲过的双键）时，第二对电子像把圆轴焊成了方轴，两端\textbf{不能}自由转动——想转就得先把第二根键拆开，那要付一大笔能量。双键两端的原子一旦被固定在同侧或异侧，就锁死了。
\end{itemize}

“会转”和“锁死”的差别听起来小，后面你会看到它撬动了多大的世界：塑料能不能拉成丝（第 43 章）、脂肪是油还是蜡（第 40 章）、乃至细胞膜是软是硬，都归这两条关节规则管。

\section{把碳骨架画到纸上：四种画法}

分子看不见，化学家得把它们画到纸上。几百年下来形成了四种“画语”，精度层层递进。以四个碳的直链分子丁烷为例：

\begin{itemize}[itemsep=2pt]
\item \textbf{分子式}：\chem{C_4H_{10}}。只报“原料清单”——几个碳几个氢，不告诉你怎么连。信息最少，写信封够用。
\item \textbf{结构简式}：\chem{CH_3-CH_2-CH_2-CH_3}。把每个碳抓了几个氢写出来，用短横表示碳碳之间的连接。怎么连一目了然，但仍是平面的。
\item \textbf{骨架式}（键线式）：更省事的速写——只画碳碳之间的键，连成折线；\textbf{每个端点和拐点都是一个碳}，碳字不写出来，碳上抓的氢也全省略（反正每个碳凑够四只手，缺几只补几个氢）。有机化学家日常聊天用的就是这种画语，本书后面越用越多。
\item \textbf{球棍模型}：小球代表原子，小棍代表共价键，搭出立体的四面体和键角。信息最全，但也最占地方。
\end{itemize}

\begin{tikzpicture}[scale=1.1]
\draw[thick] (0,0) -- (1,0.5) -- (2,0) -- (3,0.5);
\fill (0,0) circle (2pt);
\fill (1,0.5) circle (2pt);
\fill (2,0) circle (2pt);
\fill (3,0.5) circle (2pt);
\node[below] at (1.5,-0.15) {\small 丁烷的骨架式：四个点代表四个碳，氢全部省略};
\end{tikzpicture}

请务必记住：这四种画法全是\textbf{人为约定的表示法}。真实的分子里没有小棍，电子也不是两颗待在棍上的珠子（第 19 章已经强调过）；骨架式里消失的碳和氢更不是不存在，只是“读者默认知道”。画语越简洁，需要你在脑子里补的信息越多——这正是读有机结构图的基本功。

\section{化学家怎么知道碳是四面体的}

你说，分子那么小，谁见过四面体？问得好。这正是本章最值得讲的侦探故事。

\begin{zhengju}
19 世纪中叶，化学家已经确认碳是四价的，但没人知道这四根键在空间里朝哪。1874 年，22 岁的荷兰青年范特霍夫和 27 岁的法国青年勒贝尔，几乎同时、各自独立地发表了同一个大胆主张：\textbf{碳的四根键指向四面体的四个顶点}。

他们的推理漂亮在：不用看见分子，只用\textbf{数异构体的个数}。先看事实：把甲烷 \chem{CH_4} 里的一个氢换成氯，得到的 \chem{CH_3Cl} 全世界只有一种；换成两个氯，得到的 \chem{CH_2Cl_2} 化学家反复制备、精馏、测熔点沸点，几十年来也只找得到\textbf{一种}。现在做假设检验：假如四根键在同一个平面上（比如正方形四角），那么两个氯摆在相邻角和摆在对角，应该是\textbf{两种}不同的 \chem{CH_2Cl_2}——性质应有差别，应该能分开。可实验死活只给出一种。平面假设被证据逼到墙角。而四面体假设下，任意两个顶点都相邻，二氯甲烷天然只有一种——预言与事实吻合。

替代解释有没有被认真排除？有的。怀疑者说：也许第二种 \chem{CH_2Cl_2} 只是还没人造出来？于是化学家用各种路线反复合成了几十年，一无所获。又有人说：也许两种确实存在，但互相转化太快、分不开？这个反驳很有水平，它逼出了更强的证据——四面体模型预言，一个碳连着四个\textbf{各不相同}的基团时，应该存在一对“左手版和右手版”的分子，物理性质几乎全同，却互为镜像、无法重合，而且绝不可能靠转动互变。这样的分子后来真的被成对找到了（早在 1848 年，巴斯德就用镊子在显微镜下手分过酒石酸的两种晶体，当时没人说得清为什么）。范特霍夫的小册子刚出版时被大牌化学家公开嘲讽，说他是“骑在飞马上的空想家”；但 20 世纪初，布拉格父子用 X 射线衍射直接测出金刚石中每个碳的四个邻居确实占据四面体顶点——空想落了地。范特霍夫后来拿到了历史上第一届诺贝尔化学奖。

这个故事的方法论和门捷列夫如出一辙：从“数得清的个数”出发，让所有错误的形状假设无处藏身，再用可被推翻的预言（镜像分子必须成对存在）押注。你不需要看见分子，也能知道分子的形状——只要你数的数足够硬。
\end{zhengju}

\section{同样的原子，不同的分子：同分异构体}

现在给前面反复出现的“岔路”一个正式名字。分子式相同（原料清单完全一样）、结构却不同的分子，互为\textbf{同分异构体}。它们不是“同一物质的不同写法”，而是货真价实的不同物质：沸点不同、气味不同、毒性不同，能被分离、能各自装瓶。按“不同”的层级，异构体分三大家族：

\textbf{第一家族：构造异构——连接方式不同。}四个碳的 \chem{C_4H_{10}} 有直链和带支链两种连法，这就是正丁烷和异丁烷：正丁烷沸点约 $-0.5\,^{\circ}\mathrm{C}$，异丁烷约 $-12\,^{\circ}\mathrm{C}$，打火机和喷雾罐里装的往往是它们的混合物。连接方式不同还包括“同一个挂件挂在链上不同位置”“同一个分子式拆成不同种类的挂件”——这两类都和挂件本身有关，第 37 章讲官能团时会细说。

\textbf{第二家族：顺反异构——锁死的关节两侧摆位不同。}双键不能自由转动，于是双键两端的基团摆在一侧（顺式）还是分处两侧（反式），就是两种分子。比如两个碳用双键相连、各带一个甲基，顺式沸点约 $4\,^{\circ}\mathrm{C}$，反式约 $1\,^{\circ}\mathrm{C}$——差别不大，但足够分瓶。这个差别在生命世界里会被放大：天然植物油里的双键大多是顺式，让油脂分子弯成一个“肘”、堆不严实，所以常温是液体；部分加工过程会把一些双键掰成反式，分子变直、堆得紧密，熔点升高——这就是配料表上“反式脂肪酸”的由来，过量摄入与心血管风险相关（剂量和证据的问题，第 49 章讲脂类时再展开）。

\textbf{第三家族：对映异构——互为镜像、无法重合。}一个碳连着四个各不相同的基团时，这个碳就像一只手：分子有左手版和右手版，分子式相同、连接顺序也相同，却是两种物质——范特霍夫预言的那一对。这对“左右手分子”的脾气差异能大到离谱，因为酶和受体本身也是“有手性”的，认左手不认右手。这背后的故事太长也太重要，第 42 章专门讲，这里先记住一句话：\textbf{镜像，在分子世界里不总是等价的}。

一个分子可以同时叠好几层异构：既能和别的分子比连接方式，自己内部又可能有双键摆位、有镜像。碳骨架越搭越大，这些“同分不同构”的岔路就呈指数铺开——这正是有机分子百万大军的兵源。

\section{这些模型的边界在哪里}

老规矩，模型有使用说明书，也有失效区。

第一，\textbf{所有画法都是表示法}。球棍模型把原子画成硬球、键画成小棍，方便你想三维关系，但真实的原子和键不是这么回事；骨架式更是“满纸皆省略”。用它们推理可以，别把它们当照片。

第二，\textbf{“单键自由旋转”是近似}。转动其实要翻过一道小小的能量坎，低温下会转得犹豫；而一旦碳进了环，旋转被整个锁死，环的几何形状就变得非常重要。另外，三个碳、四个碳咬成的小环特别“憋屈”：四面体想要 $109.5^{\circ}$，环丙烷硬把键角掰到 $60^{\circ}$，键被掰弯、能量升高，这种“环张力”让小环分子格外容易开环反应。本章的“随便连成环”对三元环、四元环要打折扣。

第三，\textbf{“数异构体”的账我们故意算小了}。前面那张表只数了连接方式；若再把顺反、镜像都算上，种数还要多得多。反过来，账也可能算大：有些纸上画得出的结构因原子互相“撞车”（空间位阻）而极不稳定，现实中未必造得出来。模型给的是可能性清单，大自然和实验室再做第二轮筛选。

第四，\textbf{有一种碳骨架本章的画法根本画不像}：苯环。它的六个碳上没有真正意义上的单键双键之分，电子是“匀”在整个环上的——强行用单双键交替去画，会漏掉它出奇稳定的本性。这是第 41 章的主角，这里只立个警示牌。

\begin{wuqu}
“有机物就是天然物，来自生命；天然的就是安全健康的，人工合成的都有害。”——两句话，双双错误，而且错误的方式值得拆开看。

先看“有机”这个历史包袱。19 世纪初，化学家确实相信有机物只能由生命体内的“生命力”制造，实验室造不出来。1828 年，德国化学家维勒本想加热一种无机盐（氰酸铵），却意外得到了尿素——一种货真价实的、哺乳动物尿里的有机物。烧杯里没有生命力，尿素照样生成。“有机=生命制造”的墙就此塌了：今天“有机化学”只剩一个技术定义——碳骨架的化学，跟来源无关。实验室能合成胰岛素，细菌也能“天然生产”肉毒毒素。

再看“天然=安全”。毒蘑菇里的鹅膏毒素、河豚体内的河豚毒素、发霉花生上的黄曲霉毒素，全是纯天然的有机分子，个个剧毒；而阿司匹林、胰岛素、疫苗佐剂是人工合成或半合成的，救了无数人的命。判断一个物质安不安全，看的是它的结构、剂量、接触途径和证据（第 35 章的侦探态度），不是看它的“出身”。顺便说，护肤品广告里“不含化学成分”是一句语义上的空话——水都是化学物质，你自己就是一座化工厂。
\end{wuqu}

\begin{shiyan}
\textbf{动手搭一搭碳骨架}（在家可做，材料安全）

材料：橡皮泥两三种颜色（或用去壳熟鹌鹑蛋、葡萄、棉花糖代替），一盒牙签，一把尺子。

步骤一：搭甲烷。搓一个大球当碳、四个小球当氢，把四根牙签插进大球、另一端插上小球。调整四根牙签，让\textbf{任意两根的夹角尽量相等}——你会发现，在平面上怎么摆都有一对夹角偏小，只有把一根竖起来、让四根指向空间四个方向，角度才均匀。用尺子量一量相邻牙签尖端的距离：四面体摆法下，六个距离应大致相等。你亲手摸到了 $109.5^{\circ}$。

步骤二：搭两个丁烷。取四个中球当碳，先用三根牙签连成一条直链，再拆开来，改成三个连成链、一个挂在中间碳上。数一数：球和牙签的数量变了吗？没变——但形状再也对不上了。这就是 \chem{C_4H_{10}} 的两种同分异构体。

步骤三：把六个碳的链首尾相咬，搭一个六元环。试着让每根牙签的夹角接近四面体角——你会发现六元环自然翘成一把“小椅子”，并不平躺。

步骤四：搭一个“锁死的关节”。两个球之间并排插\textbf{两根}牙签当双键，两端各插上小球。现在试着像转动单键那样扭它——牙签会别住，要么扭不动，要么折断。双键为什么不能自由旋转，你的手已经知道答案了。

观察点：步骤一里“平面摆不均匀、立体才均匀”；步骤二里“同样多的球、对不上的形状”；步骤四里牙签被别住的手感。

安全提示：牙签尖锐，别对着眼睛和脸，也别含在嘴里；如果用了葡萄或棉花糖，搭完请洗手，沾过橡皮泥和牙签的食物\textbf{不要吃}。
\end{shiyan}

\section{回到开场的四样东西}

现在可以对答案了。铅笔芯和钻石为什么天差地别？第 24 章已经拆过这个案子：钻石里每个碳都用四只手拉住四个邻居，织成一张伸向四面八方的立体网，想撬动任何一个碳都得同时折断好几根键，所以硬到极致；而石墨里每个碳只伸出三只手，连成一张张平整的六边形“渔网”，层与层之间没有真正的化学键，只靠微弱的分子间作用力（第 22 章）叠着，一推就滑——你写字，就是石墨一层层滑脱、留在纸上的过程。同一种原子，\textbf{连接方式不同}，于是一个是划玻璃的王，一个是写字的笔。

糖和塑料袋呢？它们的碳骨架上挂着不同的“挂件”：糖挂满亲水的小爪子，所以溶于水；聚乙烯（第 43 章）就是一条纯粹的碳链，挂上成千上万个氢，链长到能缠成一锅原子级的面条，所以柔韧不透水。你的手背——皮肤里的蛋白质、脂肪、DNA，全是不离四只手规则的碳骨架，只是挂件的种类和排列精巧到另一个量级。挂件是什么、怎么决定分子的性格，马上就来。

你猜到了吗？答案就是贯穿问题第②问的全部威力：组成一样，\textbf{连接和排列}不同，世界就完全不同。第一册教你看“由什么组成”和“变不变、快不快”；从这一章起，“怎么连、怎么排、朝哪边”成为主角。

\section{骨架搭好了，接口呢}

这一章我们只关心骨架：碳链多长、怎么分叉、有没有环、关节锁不锁。可光有骨架的分子是“哑巴”——纯碳链能做的事很有限，无非是烧掉（第 39 章）或者缠成塑料（第 43 章）。真正让一个有机分子会酸、会香、会醉人、会治病的，是骨架上那些特定的“接口”：一个氧挂在这里，分子就变成了酒；多挂一个氧，就变成了醋；换个位置挂氮，又可能变成某种药。

这些接口有统一的名字，叫\textbf{官能团}。下一章，我们就去认识这套接口的“标准件家族”——你会发现，百万种有机分子的脾气，其实由区区十几种接口决定。

\begin{tiaozhan}
\textbf{挑战：你能不重不漏地数出 \chem{C_6H_{14}} 的全部 5 种异构体吗？}（跳过本节不影响主线。）

乱画必重必漏，需要一套“系统搜索”策略。诀窍是从最长的主链开始，逐步缩短：

第一步，主链 6 个碳：只有一种（一条直链）。

第二步，主链 5 个碳，剩下 1 个碳当支链：支链只能挂在主链第 2 位或第 3 位（挂第 1 位等于主链变长，挂第 4 位与第 2 位对称重复），得 2 种。

第三步，主链 4 个碳，剩下 2 个碳：可以两个都当甲基，也可以合成一个乙基——但乙基挂上去会让主链变相变长，与前面的重复，所以只剩“两个甲基”的情形：挂 (2,2) 位或 (2,3) 位，得 2 种。

第四步，主链 3 个碳：装不下 3 个碳的支链而不重复，0 种。

合计 $1+2+2=5$ 种，与正文表格吻合。这套“缩短主链、移动支链、随时检查重复”的流程，其实就是计算机枚举分子的算法的直觉版。你可以拿 \chem{C_7H_{16}}（答案：9 种）练手，并体会一下：为什么人类干脆写了程序来数它们？
\end{tiaozhan}

\section*{练一练}

\begin{sikaoti}
\item 不看任何资料，写出 \chem{C_5H_{12}} 的全部 3 种同分异构体的结构简式，并各画一个骨架式。检查一遍：每个碳是否恰好四只手？
\item 某同学画了两种“不同”的丁烷结构，第一种是直链、第二种把直链中间折了个弯。它们是同分异构体吗？用“单键是会转的关节”这一规则说明理由。
\item 乙烯 \chem{C_2H_4} 的两个碳之间是双键。把每个碳上的两个氢各换掉一个、换成氯，一共能得到几种不同的分子？如果换成乙烷 \chem{C_2H_6}（单键）做同样的替换，种数为什么不一样？请画出粒子摆放示意图并说明。
\item 金刚石和石墨都是纯碳。用本章的“连接和排列”语言，向一个没学过化学的朋友解释：为什么钻石能划开玻璃，铅笔芯却一写就掉粉？（不许用“硬度不同所以硬度不同”式的循环解释。）
\item 骨架式里，一个端点和一个拐点各代表什么？请把结构简式 \chem{CH_3-CH(CH_3)-CH_2-CH_3} 改画成骨架式，并说出这个分子里有几个碳、几个氢。
\item 范特霍夫时代没有 X 射线，他却断言碳的四根键指向四面体。他用的是什么类型的证据？请模仿他的思路，设计一个“数个数”的推理：如果碳的四根键中三根平躺在纸面、一根垂直纸面（“金字塔形”），\chem{CH_2Cl_2} 应该有几种？这个模型是被什么事实淘汰的？
\end{sikaoti}
